\item Demuestre que la variación de presión atmosférica con la altitud está dada por $P = P_0 e^{-\alpha y}$, donde $\alpha = \rho_0 g / P_0$, $P_0$ es la presión atmosférica en algún nivel de referencia $y = 0$ y $\rho_0$ es la densidad atmosférica a este nivel. Suponga que la disminución en presión atmosférica sobre un cambio infinitesimal en altura (de modo que la densidad es aproximadamente uniforme) puede expresarse, de la ecuación 14.4, como $dP = -\rho g dy$. También suponga que la densidad de aire es proporcional a la presión, lo cual, como se verá en el capítulo 18, es equivalente a suponer que la temperatura del aire es la misma para toda altitud.
